\section{Editorial method, attribution, and mathematical status}

\subsection{What this edition is}

The primary sources are the mathematical works written by Maxwell Charles
Siegel and distributed through arXiv or through his author-controlled site.

The purpose of the present text is narrower and more exact.  It fixes a
canonical witness for each work, states how older versions relate to it,
places any running notation beside the typed source notation through an
explicit crosswalk, and arranges the results by mathematical dependence rather
than by upload date.  The result is a modern recasting of the corpus: neither
a facsimile anthology nor a survey that replaces formulas by descriptions.

\subsection{Four statement statuses}

The edition uses the following four dispositions.  A source-issue or status
environment states its disposition locally; Chapter~10 records the source
identity and disposition of each theorem-like item admitted to its result
table.  A theorem-style heading by itself neither assigns authorship nor
certifies that a source proof has been completed.

\begin{enumerate}[label=\textup{(\roman*)}]
\item A \emph{reconstructed theorem} is a statement in Siegel's corpus whose
      source hypotheses and conclusion are preserved and placed beside any
      running notation through an explicit typed crosswalk, and whose proof is
      reproduced or reconstructed at the level needed by the subsequent
      chain.
\item A \emph{source theorem} is stated as a theorem of the named source.  If
      this edition gives only its proof mechanism, that limitation is said
      explicitly; the label does not mean that this edition supplies an
      independent proof of every auxiliary assertion in the source.
\item A \emph{source issue} records a typing error, a changed convention, an
      omitted hypothesis, a withdrawn paper, or a point at which the source
      itself supplies only a sketch.  The issue is not repaired by pretending
      that the altered statement was what the source always said.
\item An \emph{editorial proposition} is a formal consequence proved in this
      edition in order to join two source formulations.  It is marked as
      editorial and is not attributed to Siegel as an original theorem.
\end{enumerate}

Extraction is not verification.  Conversely, a typographical defect does not
by itself refute the intended theorem.  The edition therefore distinguishes
the exact source string, the mathematically typed reconstruction, and the
status of the proof.

\subsection{Notation policy}

The source convention $x\overset{a}{\equiv}y$ is rendered as
$x\equiv y\pmod a$.  The source symbol $[x]_a$ is retained for the standard
representative when it is part of a digit formula.  Iverson brackets are
written $\ind{P}$.  Composition is read from right to left.  Thus, for a word
$\mathbf j=(j_0,\ldots,j_{m-1})$, the ordered branch composition used below
is
\[
 H_{\mathbf j}=H_{j_0}\circ H_{j_1}\circ\cdots\circ H_{j_{m-1}}.
\]
The source sometimes indexes a word from $1$ and sometimes from $0$; the
edition uses $0$-based digit indices and states the convention whenever an
order-sensitive identity is used.

The name \emph{$(p,q)$-adic analysis} is retained as Siegel's name for
analysis of functions whose domain carries a $p$-adic topology and whose
values lie in a field with a distinct $q$-adic (or, more generally,
place-dependent) topology.  This is not the standard adjective ``adic''
applied to a single topological ring; the domain and codomain places must
always be displayed.

The name \emph{rising-continuity} is retained for the source condition
\[
 f(z)=\lim_{n\to\infty} f([z]_{p^n}),
\]
where the limit is taken in the specified codomain.  In standard language it
is continuity along the canonical truncation sequence, not automatically
topological continuity with respect to every sequence converging to $z$.
The distinction is used in \cref{sec:rising-fseries}.

\subsection{Version precedence}

The latest non-withdrawn arXiv version controls a work unless an
author-controlled later witness explicitly corrects it.  A later technical
manual controls a definition when it declares that earlier definitions are
inconsistent and fixes a forward convention.  Under this rule the 2026
Hydra/Numen manual controls the basic definitions of Hydra maps, numens, and
the general Correspondence Principle; the 2022 dissertation remains the
largest source for the analytic development.  A withdrawn record remains in
the historical corpus but cannot control a theorem contradicted by its own
withdrawal notice.

\begin{sourceissue}[Canonical does not mean infallible]
The word \emph{canonical} refers to source identity, version precedence, and
the stable presentation adopted by this edition.  It does not convert an
author claim into a proved theorem.  In particular, the 2026 paper labels its
Correspondence Principle proof a sketch; this edition preserves that fact and
separates the fully checked affine-word calculation from the remaining
dynamical implications.
\end{sourceissue}

\subsection{The exactness standard}

Each comparison is attached to an explicit object.  For example, the chain
used for a characteristic distribution is the following typed sequence, not
the sentence ``Fourier analysis detects periodic points.''  Starting from a
centered admissible pre-Hydra
\[
 \mathcal H=(K,\Lambda,(H_j)_{j\in\mathcal A_p})
 \in\{(K,\Lambda,\mathbf H):H_j\in\Aff_{\F}(K)\},
\]
affine word multiplication gives
\[
 (M_H,X_H)\in
 \operatorname{Map}(\Sigma_p^*,\operatorname{GL}_{\F}(K))
 \times\operatorname{Map}(\Sigma_p^*,K).
\]
Centeredness proves that the second component descends along
$\Dig:\Sigma_p^*\twoheadrightarrow\N_0$ to
$X_H:\N_0\to K$.  At a place $\ell$ let
\[
 D_{H,\ell}=\{z\in\Z_p^{\mathrm{dig}}:
             (X_H([z]_{p^N}))_{N\geq0}\text{ converges in }K_\ell\}.
\]
The truncation limit is a map
\[
 X_{H,\ell}:D_{H,\ell}\longrightarrow K_\ell.
\]
These measurability assertions are automatic.  Each truncation map
$S_N(z)=X_H([z]_{p^N})$ is locally constant.  Writing
$d_\ell(x,y)=|x-y|_\ell$, completeness of $K_\ell$ gives
\[
 D_{H,\ell}
 =\bigcap_{r\geq1}\bigcup_{N_0\geq0}
   \bigcap_{m,n\geq N_0}
   \{z:d_\ell(S_m(z),S_n(z))<1/r\},
\]
so $D_{H,\ell}$ is Borel.  On that Borel subspace,
$X_{H,\ell}=\lim_NS_N$ is Borel as the pointwise limit of metric-space-valued
Borel maps.  Thus only conullness remains a hypothesis.  If
$m_p(D_{H,\ell})=1$, then the pushforward is the Borel probability
\[
 \nu_{H,\ell}=(X_{H,\ell})_*
   (m_p|_{D_{H,\ell}})\in\operatorname{Prob}(K_\ell).
\]
Its Fourier--Stieltjes transform is the bounded continuous function
\[
 \widehat\nu_{H,\ell}:\widehat{K_\ell}\longrightarrow\C,
 \qquad
 \widehat\nu_{H,\ell}(\gamma)
 =\int_{K_\ell}\gamma(-x)\,d\nu_{H,\ell}(x).
\]
If a cyclic-translate space is subsequently considered, translations are
indexed by the character group, not by $K_\ell$: for
$\eta\in\widehat{K_\ell}$,
$(\tau_\eta f)(\gamma)=f(\gamma+\eta)$, and the asserted closure must name
its ambient function space and norm.  No self-duality of $K_\ell$ is being
silently chosen here.  Likewise, the numen is not merely ``a function
encoding the dynamics'': it is the translation term of an ordered affine
branch composition and satisfies a digit recursion with a specified initial
condition.  This level of explicitness is the governing editorial standard
throughout the edition.
